% GENERATED FILE. Edit canonical lessons, then run npm run generate:books.
% canonical-source-hash: fnv1a64:d01b6f0783ac3606
% canonical-lessons: KA-C357-vancisu, KA-C357-gaupya, KA-C357-muccidu, KA-C357-bahirangapadisu, KA-C357-sancu

\chapter{A Hidden Fraud}
\label{ch:ka-a-hidden-fraud}
\hlchaptermodality{357}

\begin{chapteropening}
\textbf{By the end of this chapter:} \emph{I can say to swindle, confidential, to cover up, to reveal and a plot.}

Complete the last lesson of chapter 357: I can say to swindle, confidential, to cover up, to reveal and a plot.
\end{chapteropening}

\section[\texorpdfstring{\kn{ವಂಚಿಸು} (vañcisu)}{ವಂಚಿಸು (vañcisu)}]{\kn{ವಂಚಿಸು} (vañcisu) — to swindle; to cheat}

\label{lesson:KA-C357-vancisu}



\begin{quote}
Before the new one: say the Kannada for a ballad, then the Kannada for a moral.
\end{quote}

\subsection*{You'll want to know: \kn{ವಂಚಿಸು}}
\textbf{\kn{ವಂಚಿಸು}} — \emph{vañcisu} — \textquotedblleft{}to swindle; to cheat\textquotedblright{}.

\begin{quote}
\kn{ಆ} \kn{ಮನುಷ್ಯ} \kn{ಊರಿನ} \kn{ಜನರನ್ನು} \kn{ವಂಚಿಸಿ} \kn{ಓಡಿಹೋದನು}.
\end{quote}

\begin{quote}
\emph{ā manuṣya ūrina janarannu vañcisi ōḍihōdanu.}
\end{quote}

\begin{quote}
That man cheated the people of the town and ran away.
\end{quote}

Past (I …): \kn{ವಂಚಿಸಿದೆ}, \emph{vañciside}.

The person cheated takes -\kn{ಅನ್ನು}. \kn{ಮೋಸ} \kn{ಮಾಡು} says the same in everyday talk.

\begin{cousinweb}
\textbf{\kn{ವಂಚಿಸು}} (\emph{vañcisu}) is built on Sanskrit \emph{vañcana}, \textquotedblleft{}cheating, deception, fraud\textquotedblright{} (Monier-Williams), with the Kannada verb ending -\emph{isu}.
\end{cousinweb}

\begin{grammarlens}[title={What you've built}]
One more word for telling what happened.
\end{grammarlens}

\subsection*{Your turn}
\begin{itemize}
\raggedright
  \item \emph{Say it:} \emph{vañcisu}
  \item \emph{Say it:} \emph{vañcisu}, once more
  \item \emph{Say it:} \emph{ā manuṣya ūrina janarannu vañcisi ōḍihōdanu.} — That man cheated the people of the town and ran away.
  \item \emph{Recall:} say the Kannada for a ballad, then the Kannada for a moral, then say \emph{vañcisu} again
\end{itemize}

\begin{tcolorbox}[breakable,colback=teal!4,colframe=teal!35!black,title={Before you move on}]
What does \kn{ವಂಚಿಸು} mean? (\textquotedblleft{}To swindle; to cheat\textquotedblright{}.) Say it once more.
\end{tcolorbox}
\section[\texorpdfstring{\kn{ಗೌಪ್ಯ} (gaupya)}{ಗೌಪ್ಯ (gaupya)}]{\kn{ಗೌಪ್ಯ} (gaupya) — confidential; secret}

\label{lesson:KA-C357-gaupya}



\begin{quote}
Before the new one: say the Kannada for a moral, then the Kannada for to swindle.
\end{quote}

\subsection*{You'll want to know: \kn{ಗೌಪ್ಯ}}
\textbf{\kn{ಗೌಪ್ಯ}} — \emph{gaupya} — \textquotedblleft{}confidential; secret\textquotedblright{}.

\begin{quote}
\kn{ಈ} \kn{ಪತ್ರ} \kn{ಗೌಪ್ಯ}, \kn{ಅದನ್ನು} \kn{ಬೇರೆ} \kn{ಯಾರೂ} \kn{ಓದಬಾರದು}.
\end{quote}

\begin{quote}
\emph{ī patra gaupya, adannu bēre yārū ōdabāradu.}
\end{quote}

\begin{quote}
This letter is confidential; nobody else may read it.
\end{quote}

\kn{ಗೌಪ್ಯವಾಗಿ} \kn{ಇಡು}, to keep confidential. \kn{ಗೌಪ್ಯತೆ} is privacy.

\begin{cousinweb}
\textbf{\kn{ಗೌಪ್ಯ}} (\emph{gaupya}) is usually traced to Sanskrit \emph{gopya}, \textquotedblleft{}to be kept secret or hidden\textquotedblright{} (Monier-Williams), from the root \emph{gup}, \textquotedblleft{}to guard, to hide\textquotedblright{}; the dictionary has no form with \emph{au}.
\end{cousinweb}

\begin{grammarlens}[title={What you've built}]
One more word for telling what happened.
\end{grammarlens}

\subsection*{Your turn}
\begin{itemize}
\raggedright
  \item \emph{Say it:} \emph{gaupya}
  \item \emph{Say it:} \emph{gaupya}, once more
  \item \emph{Say it:} \emph{ī patra gaupya, adannu bēre yārū ōdabāradu.} — This letter is confidential; nobody else may read it.
  \item \emph{Recall:} say the Kannada for a moral, then the Kannada for to swindle, then say \emph{gaupya} again
\end{itemize}

\begin{tcolorbox}[breakable,colback=teal!4,colframe=teal!35!black,title={Before you move on}]
What does \kn{ಗೌಪ್ಯ} mean? (\textquotedblleft{}Confidential; secret\textquotedblright{}.) Say it once more.
\end{tcolorbox}
\section[\texorpdfstring{\kn{ಮುಚ್ಚಿಡು} (mucciḍu)}{ಮುಚ್ಚಿಡು (mucciḍu)}]{\kn{ಮುಚ್ಚಿಡು} (mucciḍu) — to cover up; to hide}

\label{lesson:KA-C357-muccidu}



\begin{quote}
Before the new one: say the Kannada for to swindle, then the Kannada for confidential.
\end{quote}

\subsection*{You'll want to know: \kn{ಮುಚ್ಚಿಡು}}
\textbf{\kn{ಮುಚ್ಚಿಡು}} — \emph{mucciḍu} — \textquotedblleft{}to cover up; to hide\textquotedblright{}.

\begin{quote}
\kn{ಅಧಿಕಾರಿ} \kn{ತನ್ನ} \kn{ತಪ್ಪನ್ನು} \kn{ಮುಚ್ಚಿಟ್ಟನು}.
\end{quote}

\begin{quote}
\emph{adhikāri tanna tappannu mucciṭṭanu.}
\end{quote}

\begin{quote}
The officer covered up his mistake.
\end{quote}

Past (I …): \kn{ಮುಚ್ಚಿಟ್ಟೆ}, \emph{mucciṭṭe}.

\kn{ಮುಚ್ಚಿ}, having closed, and \kn{ಇಡು}, to keep: to keep closed. The past is \kn{ಮುಚ್ಚಿಟ್ಟೆ}.

\begin{cousinweb}
\textbf{\kn{ಮುಚ್ಚಿಡು}} (\emph{mucciḍu}) is a native Dravidian compound: \emph{muccu}, \textquotedblleft{}to close, to cover, to conceal\textquotedblright{} (DEDR 4915), and \emph{iḍu}, \textquotedblleft{}to put, to keep\textquotedblright{} (DEDR 442).
\end{cousinweb}

\begin{grammarlens}[title={What you've built}]
One more word for telling what happened.
\end{grammarlens}

\subsection*{Your turn}
\begin{itemize}
\raggedright
  \item \emph{Say it:} \emph{mucciḍu}
  \item \emph{Say it:} \emph{mucciḍu}, once more
  \item \emph{Say it:} \emph{adhikāri tanna tappannu mucciṭṭanu.} — The officer covered up his mistake.
  \item \emph{Recall:} say the Kannada for to swindle, then the Kannada for confidential, then say \emph{mucciḍu} again
\end{itemize}

\begin{tcolorbox}[breakable,colback=teal!4,colframe=teal!35!black,title={Before you move on}]
What does \kn{ಮುಚ್ಚಿಡು} mean? (\textquotedblleft{}To cover up; to hide\textquotedblright{}.) Say it once more.
\end{tcolorbox}
\section[\texorpdfstring{\kn{ಬಹಿರಂಗಪಡಿಸು} (bahiraṅgapaḍisu)}{ಬಹಿರಂಗಪಡಿಸು (bahiraṅgapaḍisu)}]{\kn{ಬಹಿರಂಗಪಡಿಸು} (bahiraṅgapaḍisu) — to reveal; to expose}

\label{lesson:KA-C357-bahirangapadisu}



\begin{quote}
Before the new one: say the Kannada for confidential, then the Kannada for to cover up.
\end{quote}

\subsection*{You'll want to know: \kn{ಬಹಿರಂಗಪಡಿಸು}}
\textbf{\kn{ಬಹಿರಂಗಪಡಿಸು}} — \emph{bahiraṅgapaḍisu} — \textquotedblleft{}to reveal; to expose\textquotedblright{}.

\begin{quote}
\kn{ಒಬ್ಬ} \kn{ಪತ್ರಕರ್ತ} \kn{ಆ} \kn{ಮೋಸವನ್ನು} \kn{ಬಹಿರಂಗಪಡಿಸಿದನು}.
\end{quote}

\begin{quote}
\emph{obba patrakarta ā mōsavannu bahiraṅgapaḍisidanu.}
\end{quote}

\begin{quote}
A journalist exposed that fraud.
\end{quote}

Past (I …): \kn{ಬಹಿರಂಗಪಡಿಸಿದೆ}, \emph{bahiraṅgapaḍiside}.

\kn{ಬಹಿರಂಗ}, open, public, with -\kn{ಪಡಿಸು}, to make: to make public.

\begin{cousinweb}
\textbf{\kn{ಬಹಿರಂಗಪಡಿಸು}} (\emph{bahiraṅgapaḍisu}) is built from Sanskrit \emph{bahiraṅga}, \textquotedblleft{}relating to the exterior, external\textquotedblright{} (Monier-Williams), and the Kannada verb-maker -\emph{paḍisu}, from native \emph{paḍu} (DEDR 3853).
\end{cousinweb}

\begin{grammarlens}[title={What you've built}]
One more word for telling what happened.
\end{grammarlens}

\subsection*{Your turn}
\begin{itemize}
\raggedright
  \item \emph{Say it:} \emph{bahiraṅgapaḍisu}
  \item \emph{Say it:} \emph{bahiraṅgapaḍisu}, once more
  \item \emph{Say it:} \emph{obba patrakarta ā mōsavannu bahiraṅgapaḍisidanu.} — A journalist exposed that fraud.
  \item \emph{Recall:} say the Kannada for confidential, then the Kannada for to cover up, then say \emph{bahiraṅgapaḍisu} again
\end{itemize}

\begin{tcolorbox}[breakable,colback=teal!4,colframe=teal!35!black,title={Before you move on}]
What does \kn{ಬಹಿರಂಗಪಡಿಸು} mean? (\textquotedblleft{}To reveal; to expose\textquotedblright{}.) Say it once more.
\end{tcolorbox}
\section[\texorpdfstring{\kn{ಸಂಚು} (sañcu)}{ಸಂಚು (sañcu)}]{\kn{ಸಂಚು} (sañcu) — a plot; a conspiracy}

\label{lesson:KA-C357-sancu}



\begin{quote}
Before the new one: say the Kannada for to cover up, then the Kannada for to reveal.
\end{quote}

\subsection*{You'll want to know: \kn{ಸಂಚು}}
\textbf{\kn{ಸಂಚು}} — \emph{sañcu} — \textquotedblleft{}a plot; a conspiracy\textquotedblright{}.

\begin{quote}
\kn{ಅರಮನೆಯಲ್ಲಿ} \kn{ಕೆಲವರು} \kn{ರಾಜನ} \kn{ವಿರುದ್ಧ} \kn{ಸಂಚು} \kn{ಮಾಡಿದರು}.
\end{quote}

\begin{quote}
\emph{aramaneyalli kelavaru rājana viruddha sañcu māḍidaru.}
\end{quote}

\begin{quote}
In the palace some people plotted against the king.
\end{quote}

Plural: \kn{ಸಂಚುಗಳು}, \emph{sañcugaḷu}.

A plot is \textquotedblleft{}made\textquotedblright{}: \kn{ಸಂಚು} \kn{ಮಾಡು}, to plot.

\begin{cousinweb}
\textbf{\kn{ಸಂಚು}} (\emph{sañcu}) is native Dravidian: DEDR 2293 lists Kannada \emph{sañcu}, \textquotedblleft{}an expedient, an artifice, a stratagem, an intrigue\textquotedblright{}.
\end{cousinweb}

\begin{grammarlens}[title={What you've built}]
One more word for telling what happened.
\end{grammarlens}

\subsection*{Your turn}
\begin{itemize}
\raggedright
  \item \emph{Say it:} \emph{sañcu}
  \item \emph{Say it:} \emph{sañcu}, once more
  \item \emph{Say it:} \emph{aramaneyalli kelavaru rājana viruddha sañcu māḍidaru.} — In the palace some people plotted against the king.
  \item \emph{Recall:} say the Kannada for to cover up, then the Kannada for to reveal, then say \emph{sañcu} again
\end{itemize}

\begin{tcolorbox}[breakable,colback=teal!4,colframe=teal!35!black,title={Before you move on}]
What does \kn{ಸಂಚು} mean? (\textquotedblleft{}A plot; a conspiracy\textquotedblright{}.) Say it once more.
\end{tcolorbox}
